
\textit{This chapter summarizes the project, revisits the objectives, and reflects on the overall outcome of the developed system}


The goal of the project was to evaluate whether alternative visualization methods and features could provide new value to the Norwegian Coastal Administration, in analyzing \acrshort{ais} anomaly data. The focus was on developing a prototype application that enables analysis of large-scale anomaly data sets through interactive filtering and multiple new visualization approaches. 


The final system demonstrates a locally hosted prototype allowing for analysis of \acrshort{ais} anomaly data. It allows users to view anomaly patterns through clustering, heatmap visualizations and detailed inspection of anomaly groups and individual anomalies. In addition, filtering based on \acrshort{mmsi}, anomaly type and date range were added to enable a targeted analysis of anomalies and their connection to base station and satellites. Together, these features support a high level view of where anomalies aggregate, and their signal connection to base stations on the coast alongside satellite representation. 

The project demonstrates that it is possible for third party tools and modern web technologies to prototype new features and methods in data analysis. This approach enables early evaluation of visualization concepts without requiring full scale development on these systems, allowing verification of the concepts value and reducing development risk, and unnecessary time usage. 

The project objectives were in large part achieved, fully implementing the core visualization and interaction features requested, including clustering, heatmaps, filtering and information components, alongside signal strength and individual anomaly layers. However, some planned features, such as hexagonal binning, more complex visualization of satellite orbits, and report generation, were not implemented due to time constraints and scope cuts that kept focus on the delivered system.
% SUGGESTION (criterion 5 rewards an explicit goal-by-goal verdict): consider walking the four
% Introduction objectives (Advanced Map Visualization, Data Aggregation/Clustering, Temporal Analysis,
% Efficient Querying) and stating met / partially met / not met for each. Temporal Analysis was only
% partially met (date-range filter, not forward/backward animation) -- saying so explicitly scores
% better than a general "in large part achieved".

Further improvement could be implemented, such as integration with the customers real-time \acrshort{ais} data stream, to enable monitoring of live data. More advanced temporal analysis or implementation of cut features such as hexagonal binning, can add new value to the system. The system can also be further worked on to visualize more than \acrshort{ais} data, and can be used to test concepts related to other fields and visualization methods. 

Overall, the implemented system meets the main objective of creating a simple, readable visualization program for \acrshort{ais} anomaly data, while remaining open for further work. 



